\scanpage{112}
% archive foliation: 111

\deleted{d) Si $F$ spécial de \unsure{poids}}

\uline{Proposition} \unsure{3.3}.

a) Stabilité de ``eff. adm. de poids $\subset\Phi$'' par sous-faisceaux constructibles, quotients constructibles, \inserted{sommes finies,} extensions.

b) Si $F,F'$ sont eff adm à poids $\subset\Phi,\Phi'$, alors $F\otimes F'$ est admissible à poids $\subset\Phi+\Phi'=\{i+i'\mid i\in\Phi,\ i'\in\Phi'\}$.

c) Stabilité par \unsure{chgt} de base

d) Si $F$ est \deleted{spécial}\inserted{eff$^{\unsure{t|!}}$} admissible \uline{constant tordu} à poids $\subset\Phi$, \unsure{et} si $i\in\Phi\Rightarrow i\leq n$, alors $\check F(\unsure{-n})$ est \inserted{eff$^{\unsure{t|!}}$} admissible (constant tordu) à poids $\subset2n-\Phi=\{n-i\mid i\in\Phi\}$.

\margin{d) on admet Lefschetz}

\uline{Dém} a) b) c) sont immédiats, en utilisant les assertions correspondantes de 2.1. Notons d'ailleurs:

3.2. \inserted{Supposons $S$ qu. cpt qu. sép.} Pour que $F$ soit eff$^{\unsure{t|!}}$ adm. à poids $\subset\Phi$, il f et s. que l'on puisse trouver une partition de $S$ en \deleted{\unsure{finaux}} sous-schémas $S_i\xrightarrow{\alpha_i}S$, $\alpha_i$ immers. de prés. finie, avec \deleted{\unsure{des}} $F|S_i=F_i$ admettant une filtration dont les quotients successifs sont spéciaux de poids $\in\Phi$.

\margin{\figure{\begin{tikzpicture}[x=1cm,y=1cm]
\draw[->] (0,0) .. controls (-.7,.05) and (-.4,1.7) .. (-.4,2.1) .. controls (-.5,3.8) and (-1.3,5.2) .. (-.55,6.7);
\end{tikzpicture}}}

Comme, pour $F$ constant tordu, la formation de $\check F$ commute au changement de base, \inserted{3.\unsure{3|2}} d) résulte de 2.2.

\deleted{\unsure{ptt}}

4.1 Un faisceau $F$ est dit eff$^{t}$-admissible de niveaux $\leq\unsure{\nu|\gamma}$ si \deleted{il est effectivement admissible, et} dans 3.1, on peut prendre $G$ \deleted{de poids $j\in\unsure{\Phi}$} de la forme $G'(-k)$, $k\geq0$, où $G'$ est spécial de poids \deleted{$\unsure{j\in\Phi}$} $\leq\unsure{\nu|\gamma}$. --- On dit que $F$ est eff$^{t}$ admissible ``de poids $\subset\Phi$, et de niveaux $\leq\nu$'' si on peut choisir dans 3.1. $G'$ \deleted{\unsure{de}} spécial de poids $j\leq\nu$, avec \deleted{poids} $j+2k\in\Phi$.

\margin{\unsure{bloquer après 3.2, en énoncé analogue: 3.2 pour soit 3.1 ou} \deleted{\unsure{donc}} \unsure{que} 4.1.}

\uline{Attention}, à priori cela est plus fort que la \uline{conjonction} de ``de poids $\subset\Phi$'' et ``de niveaux $\leq\nu$''. \deleted{\unsure{Pour cela} \illegible{2} \unsure{poids}, \unsure{on a}} \inserted{On notera que si} \deleted{$\unsure{\Phi\leq d}$} $\nu$ est \deleted{\unsure{un entier}} un entier tel que $\Phi\leq\nu$ (i.e $j\in\Phi\Rightarrow j\leq\nu$), \deleted{$\unsure{\frac{n}{2}}$ \unsure{est} $\unsure{j\in\Phi}$} alors ``eff$^{\unsure{t|!}}$ adm. de poids $\in\Phi$'' $\Rightarrow$ ``id et à niveaux $\leq\nu$''.
